\section{Conclusion and Future Work}
\label{sec:conclusion}

In this work, we introduced \textbf{GSC-Merge} (Grassmannian Subspace Consensus Merging), a mathematically rigorous and provably optimal framework for multi-task parameter alignment in weight space. By replacing heuristic, coordinate-wise pruning pipelines with a principled spectral consensus projection, we showed that weight-space model merging can be grounded in the elegant geometry of the Grassmannian manifold. Leveraging the Eckart-Young-Mirsky Theorem, we proved that our SVD-based projection operator provides the mathematically optimal low-rank approximation of joint downstream trajectories, effectively filtering out destructive, high-frequency, task-specific parameter interference.

Crucially, we exposed the \emph{Overfitting-Optimizer Paradox} of unconstrained few-shot validation tuning, showing that optimizing blending coefficients on tiny calibration sets leads to transductive overfitting and poor test generalization. We showed both theoretically and empirically that GSC-Merge elegantly resolves this paradox through implicit spectral regularization, restricting updates to a low-dimensional consensus subspace and substantially boosting out-of-distribution generalizability.

Our work opens several exciting avenues for future research. First, while we focused on Singular Value Decomposition (SVD), future investigations can explore alternative manifold projections, such as Stiefel manifolds or optimal transport-based barycentric coordinates, to further align representation spaces. Second, the concept of spectral consensus filtering can be extended beyond vision models to high-capacity Large Language Models (LLMs) and diffusion backbones, where parameter interference is similarly a critical bottleneck. Finally, rather than employing a globally uniform fractional rank parameter $\gamma$, an adaptive layer-wise spectral thresholding scheme based on the local singular value decay spectrum represents a promising direction to optimize parameter capacity on a layer-by-layer basis. By transitioning from discrete heuristics to continuous spectral operators, we hope to establish a more principled foundation for deep neural network parameter fusion.
